\begin{center}
  {\zihao{3}\bfseries 摘要}
\end{center}

本文围绕多区域算力基础设施的需求预测、任务调度、储能协同与算-储-电联合决策，构造由问题一向问题四逐层传递的证据链。问题一采用“季节朴素基线--共享 HistGradientBoosting 残差校正--层级一致性校准--split-conformal 区间”的预测模型，并将实际到达任务交给带容量、功率、时延和收尾边界的 CP-SAT 调度器；独立盲测系统加权 WAPE 为 $0.8934394917$，名义 $95\%$ 区间经验覆盖率为 $94.67592593\%$，实际任务完成率为 $100\%$，求解状态为 FEASIBLE，因此仅作可行调度结论。

问题二在真实到达任务上采用“全时域确定性滚动影子交换启发式”：先筛除违反硬约束的候选，再用拉格朗日平衡项比较跨区域交换，FIFO--local-first 方案作为同输入、同约束、同输出口径的基线。50,000 个任务全部覆盖，候选完成率为 $1$、SLA 违约率为 $0$；相对 FIFO，成本变化为 $-1.0149373919\%$，碳排变化为 $-1.9292972388\%$，加权平均网络时延增加 $4.18976\,\mathrm{ms}$。该方法是有界一遍启发式，不宣称全局最优。

问题三将问题二输出的任务负荷固定为输入，在六个区域分别求解 $14\times168\,\mathrm{h}+55\,\mathrm{h}$ 的滚动二进制储能 MILP，共 90 个块，覆盖 $0$--$2406\,\mathrm{h}$；与无储能新能源优先基线相比，六区域合计成本差为 $-38,899,969.87293261\,\mathrm{CNY}$，碳排差为 $-8.8505625719\,\mathrm{tCO_2}$，四类 72 小时压力探针的候选与基线审计 48/48 通过。设施负荷复算最大残差为 $5.044995509706496\times10^{-5}\,\mathrm{MW}$，完整审计为 269/270，通过舍入残差边界解释，不将滚动结果写成单体全时域最优；含同时充放电的全周期 LP 探针永久排除。

问题四固定问题二的候选与 FIFO 任务包络，在五类情景下进行六区域、72 小时顺序耦合二进制储能 MILP，并单独扫描 24 小时低新能源峰值权重。在低新能源情景中，候选相对基线的成本差为 $-1,683,494.3771166522\,\mathrm{CNY}$、碳排差为 $-189.1085063591487\,\mathrm{tCO_2}$；峰值探针选择权重 $0.0003$ 时，系统峰值由 $291.60330044$ 降至 $197.08826651333345\,\mathrm{MW}$，削峰 $94.51503392666643\,\mathrm{MW}$。该结果是固定任务包络下的顺序耦合证据和离散权重比较，不代表任务--储能联合或全时域全局最优。确定性复跑、硬约束审计、基线配对和边界排除共同构成模型检验。

\noindent\textbf{关键词：}算力需求预测；碳感知调度；滚动二进制储能 MILP；顺序耦合；峰值权衡
